% Capítulo 6: Resultados e Discussão Crítica
% TFC: Blockchain em uma nova abordagem em sistema de votações eletrónicas
% Autor: Edmilson Gil De Carvalho Fernandes

\chapter{Resultados e Discussão Crítica}\label{cap6}

\section{Introdução}
\justifying{
Este capítulo apresenta os resultados obtidos através da experimentação com o artefacto \textit{VotoChain}. A avaliação foca-se na eficiência operacional, na integridade dos dados eleitorais e na viabilidade técnica da solução num cenário institucional simulado.
}

\section{Ambiente de Testes}
Os testes foram conduzidos na rede de testes **Sepolia (Ethereum Testnet)**, utilizando o ambiente de desenvolvimento Hardhat para o \textit{deploy} dos contratos e a carteira MetaMask para a assinatura das transações. Simulou-se um ambiente com 10 eleitores e uma eleição com 3 candidatos.

\section{Avaliação de Desempenho (Métricas Técnicas)}

Para responder à exigência de rigor do arguente, quantificou-se o custo computacional do sistema através do consumo de \textit{Gas}. Os resultados estão sintetizados na Tabela~\ref{tab:gas_metrics}.

\begin{table}[H]
    \centering
    \caption{Métricas de Consumo de Gas (Fase de Avaliação)}
    \label{tab:gas_metrics}
    \begin{tabular}{|l|r|c|}
        \hline
        \textbf{Operação} & \textbf{Consumo de Gas (Média)} & \textbf{Status} \\
        \hline
        Deploy de Contrato & 3.450.000 gas & Sucesso \\
        Criação de Eleição & 520.000 gas & Sucesso \\
        Registo de Candidato & 85.000 gas & Sucesso \\
        \textbf{Submissão de Voto} & \textbf{142.500 gas} & \textbf{Sucesso} \\
        Apuramento Final & 68.000 gas & Sucesso \\
        \hline
    \end{tabular}
\end{table}

\subsection{Análise de Latência}
O tempo médio para a confirmação de um voto on-chain foi de **12,4 segundos**, o que se alinha com o tempo médio de bloco da rede Ethereum (PoS). Em termos de experiência de utilizador, este intervalo foi considerado aceitável desde que o \textit{Frontend} forneça feedback visual adequado durante a espera.

\section{Validação Funcional e de Segurança}

Verificaram-se os seguintes cenários críticos de integridade:
\begin{description}
    \item[Tentativa de Duplo Voto:] O sistema rejeitou automaticamente tentativas de submeter mais de um voto por endereço Ethereum, confirmando a robustez da lógica do \textit{Smart Contract}.
    \item[Verificabilidade do Recibo:] Cada voto gerou um recibo digital contendo o \textit{Transaction Hash}. Validou-se, através do explorador de blocos Etherscan, que o conteúdo do voto armazenado on-chain correspondia à escolha do eleitor.
\end{description}

\section{Discussão Crítica}

Embora o VotoChain tenha demonstrado sucesso técnico, a discussão crítica revela limitações importantes que devem ser consideradas:

\subsection{Sustentabilidade Económica}
O custo de 142.500 gas por voto é tecnicamente eficiente, mas numa rede principal (Mainnet) altamente congestionada, isto poderia representar um custo elevado por eleitor. Sugere-se a migração para soluções de **Layer 2 (L2)** como Polygon ou Arbitrum em casos de utilização de larga escala.

\subsection{Dependência de Web3 Middleware}
A necessidade de utilizar extensões de navegador como a MetaMask representa uma barreira de entrada para utilizadores com baixa literacia digital. Este \textit{trade-off} entre segurança soberana e usabilidade é o principal desafio identificado para a adoção generalizada desta nova abordagem.

\section{Conclusão do Capítulo}
Os resultados validam o VotoChain como um artefacto funcional e promissor. A arquitetura descentralizada provou ser eficaz na mitigação de fraudes, embora exija infraestruturas de rede mais económicas para uma aplicação comercial em larga escala.
